La conséquence conceptuelle se concentre dans l'articulation de \(K\),
de \(\alpha\) et de la section électronique. La reconstruction réversible de
\(K\) conserve des données intervenant tant dans la normalisation positionnelle
que dans l'incidence de frontière. La composition électronique utilise les
coordonnées précédemment construites et maintient son centre, ses transports
et sa réalisation dimensionnelle. Le prolongement exceptionnel conserve d'autres
relations du même état dans des représentations successives. Les connexions
entre ces objets font donc partie du résultat, au même titre que les
valeurs de leurs coordonnées.

Cette organisation détermine aussi l'ordre de la confrontation. Lorsqu'une
dérivation fixe une valeur sans la recevoir comme argument, la mesure
met à l'épreuve la conséquence obtenue. La correspondance avec une
observable physique exige, en outre, d'identifier la grandeur, la normalisation
et la réalisation dimensionnelle employées. En ce sens précis, la thèse
examine le passage d'une description paramétrée à l'explication des
relations qui déterminent ses paramètres. La portée de chaque résultat
est celle des démonstrations exposées : structure, équation et valeur
constituent des aspects inséparables de la détermination confrontée à l'expérience.

L'unicité du registre conserve deux recensements prospectifs et leurs
opérations propres (paragraphe~\ref{act21:sec:relacion-censos}).
La chaîne \(40320\to21902\to19446\to79\to1\) aboutit au
survivant hétérotypique de la face exceptionnelle ; la chaîne
\(196\cdot197\cdot196\to331\to2\to1\) aboutit à la sélection
orientée à partir des catalogues régionaux. Les populations intermédiaires
conservent le sens de chaque filtre. La composition des sorties
sélectionnées avec l'inversion entière articule détermination
prospective et récupération réversible.
